% =============================================================================
% ENSAIO ACADÊMICO — MODELO ABNT
% =============================================================================
% Características do gênero ensaio:
%   - Tom autoral e argumentativo
%   - Base teórica (não é opinião pessoal)
%   - Reflexão crítica sobre um tema
%   - Estrutura livre mas com introdução, desenvolvimento e conclusão
%   - Indicação de 5-15 referências (dependendo da profundidade)
% =============================================================================
\documentclass[
  12pt,
  a4paper,
  brazilian
]{article}

% --- Pacotes ---
\usepackage[utf8]{inputenc}
\usepackage[T1]{fontenc}
\usepackage[brazilian,portuges]{babel}
\usepackage{amsmath,amssymb}
\usepackage[alf]{abntex2cite}
\usepackage{abntex2abrev}
\usepackage{setspace}
\usepackage{geometry}
\usepackage{hyperref}
\usepackage{microtype}

% --- Margens ABNT ---
\geometry{
  a4paper,
  left=3cm,
  top=3cm,
  right=2cm,
  bottom=2cm
}

% --- Espaçamento 1,5 ---
\onehalfspacing

% --- Citação longa recuada ---
% Usar o ambiente \begin{citacao} ... \end{citacao}

% --- Metadados ---
\hypersetup{
  pdfauthor={Nome do Autor},
  pdftitle={Título do Ensaio},
  colorlinks=true,
  linkcolor=black,
  citecolor=black,
  urlcolor=black
}

\title{Título do Ensaio: Subtítulo (se houver)}
\author{Nome Completo do Autor \\ \small{Instituição} \\ \small{e-mail}}
\date{2026}

\begin{document}

\maketitle

\section{Introdução}

O ensaio acadêmico se distingue do artigo científico por seu caráter mais
ensalístico e autoral, sem exigir coleta de dados primários ou aprovação
em comitê de ética. Exige, contudo, rigor teórico, consistência
argumentativa e domínio da literatura pertinente.

A introdução deve apresentar o tema, sua relevância e a tese central que o
ensaio defenderá. Deve conter também indicação da estrutura do texto.

\section{Desenvolvimento — Primeiro Movimento Argumentativo}

Desenvolvimento do primeiro núcleo argumentativo. Fundamentar com citações
diretas e indiretas da literatura de referência.

Segundo \cite{autor2021}, a argumentação deve ser construída com base em
evidências teóricas sólidas.

\subsection{Primeiro Subtema}

Aprofundamento do aspecto específico dentro do primeiro movimento.

\subsection{Segundo Subtema}

Continuação com eventual contraponto ou nuance.

\section{Desenvolvimento — Segundo Movimento Argumentativo}

Estabelecimento da relação entre o primeiro argumento e um aspecto
complementar ou antitético.

\section{Desenvolvimento — Síntese e Proposição}

Articulação dos argumentos anteriores em direção à tese central do ensaio.
Apresentação da posição crítica do autor.

\section{Considerações Finais}

Retomada da tese, síntese dos argumentos, implicações e questões abertas
para reflexão futura. Diferentemente da conclusão de um artigo, as
considerações finais do ensaio podem ser propositivamente inconclusivas,
convidando ao debate.

\section*{Referências}

\bibliographystyle{abntex2-alf}
\bibliography{ensaio_modelo}

\end{document}
